\section{Key Features}
\label{sec:features}

\subsection{Comprehensive Method Library}

Table~\ref{tab:methods} summarizes the methods integrated in Know-Surgery across three operation categories, totaling 35 text-based methods plus 13 multimodal methods.

\begin{table}[t]
\centering
\small
\setlength{\tabcolsep}{2pt}
\begin{tabular}{@{}llp{3.8cm}@{}}
\toprule
\textbf{Op.} & \textbf{Category} & \textbf{Methods} \\
\midrule
\multirow{5}{*}{\rotatebox{90}{Edit (15)}}
 & Locate-Edit & ROME, MEMIT, MEMIT\_M, AlphaEdit \\
 & Memory & GRACE, WISE, IKE, SERAC \\
 & Meta-learn & MEND, MALMEN, InstructEdit \\
 & Unified & UniKE, NMKE \\
 & Unstruct. & UNKE, AnyEdit \\
\midrule
\multirow{4}{*}{\rotatebox{90}{Unl. (14)}}
 & Gradient & GradAscent, GradDiff, SGA \\
 & Preference & NPO, DPO, SimNPO \\
 & Represent. & RMU, FLAT, SOUL \\
 & Others & UNDIAL, CEU, SatImp, WGA, PDU \\
\midrule
Inj. (6) & PEFT & LoRA, DoRA, AdaLoRA, LoReFT, BREP \\
\midrule
\multirow{2}{*}{MM (13)}
 & MM Unl. & MMGradAsc., MMNPO, MMVKD, ... (7) \\
 & MM Edit & MM-IKE, MM-UniKE, ... (6) \\
\bottomrule
\end{tabular}
\caption{Methods supported in Know-Surgery. MEMIT\_M = MEMIT\_Merge.}
\label{tab:methods}
\end{table}

\subsection{Multi-type Knowledge Coverage}

Unlike existing tools that focus primarily on factual triplets, Know-Surgery covers five knowledge types through 8+ editing datasets:

\begin{table}[t]
\centering
\small
\setlength{\tabcolsep}{2pt}
\begin{tabular}{@{}lp{2.6cm}p{2.8cm}@{}}
\toprule
\textbf{Type} & \textbf{Datasets} & \textbf{Methods} \\
\midrule
Factual & CF, ZsRE, LEME & ROME, MEMIT, AlphaEdit \\
Event & ELKEN, UnKE, AKEW & MEMIT, AnyEdit \\
Concept & ConceptEdit & ROME, AlphaEdit \\
Proced. & EditEvery & AnyEdit, UNKE \\
Rule & RuleTaker, FOLIO & BREP, LoReFT \\
\bottomrule
\end{tabular}
\caption{Knowledge type coverage. CF=CounterFact.}
\label{tab:knowledge}
\end{table}

\subsection{Multimodal Extension}

Know-Surgery extends to multimodal scenarios through two pathways:
\begin{itemize}
\setlength\itemsep{0pt}
\item \textbf{MM Unlearning} (7 methods): MMGradAscent, MMGradDiff, MMKLMin, MMNPO, MMRetainFT, MMUnlearner, MMVKD, evaluated on CLEAR and MLLMU-Bench.
\item \textbf{MM Editing} (6 methods): MM-IKE, MM-GRACE, MM-WISE, MM-MEND, MM-SERAC, MM-UniKE, extending text editing methods to vision-language models (InternVL, LLaVA, Qwen-VL) via a shared \texttt{MMEditMixin} base class.
\end{itemize}

\subsection{Three-Step Extension Protocol}

Adding a new component requires only three steps:
\begin{enumerate}
\setlength\itemsep{0pt}
\item \textbf{Implement}: Write a Python class inheriting the appropriate base class (\texttt{EditTrainer}, \texttt{FinetuneTrainer}, or \texttt{InjectTrainer}).
\item \textbf{Register}: Add one line to the registry (\texttt{\_register\_trainer(MyMethod)}).
\item \textbf{Configure}: Create a YAML file specifying the handler name and method-specific hyperparameters.
\end{enumerate}
No changes to the core framework, data loading, or evaluation code are needed.
